\section{Twenty-four unrestricted integer unknowns}
\label{sec:integer}

The positive proof above does not extend merely by allowing \(s\)
to be negative: its comparison \(A<G\) and the signs of its weights
were proved on the positive domain. For the integer construction we
instead retain the polynomial norm circuit of package 21 and prove a
larger domain for it.

Use the common abbreviations (2.4)--(2.7), but now set
\[
\mathcal B=v^2+D+2b,\qquad
L=1+H\mathcal B,\qquad M=L,\qquad W=b_0+Hr.
\tag{6.1}
\]
Require positivity only for
\[
J,F,\alpha,z,\eta,\zeta.
\]
For each of these six coordinates use a distinct triple of integer
unknowns and the polynomial
\[
\Pos(x,y,t)=1+x^2+y^2+t(t+1).
\tag{6.2}
\]
It is positive for every integer triple and maps onto the positive
integers. Indeed \(4n-3\equiv1\pmod4\) is a sum of three squares
with one odd square and two even squares; dividing after separating
the odd square gives (6.2).

For the remaining coordinates put
\[
i=2i_0+1,\qquad s=2s_0+1,\qquad w=w_0^2,
\qquad h,T,r\in\Z.
\tag{6.3}
\]
There are \(6\cdot3+6=24\) independent integer unknowns.
After these substitutions the polynomial is
\[
\boxed{\quad P_{24}:=N_kN_t
 \bigl(1+\mathcal N(L,L,W)\bigr)-1=0.\quad}
\tag{6.4}
\]

\begin{theorem}\label{thm:integer}
For the complete compiler tuple of Section~\ref{sec:coefficients},
\(P_{24}=0\) has an integer solution if and only if \(\KTM\)
halts on its blank tape. Its exact degree is \(2432\), and its
circuit uses \(163\) binary operations.
\end{theorem}

\subsection{Bounds valid for both signs of the scale}

The unit guard again gives \(\mathcal N=0\) and
\(N_k=N_t=\epsilon\in\{1,-1\}\). Positivity of \(F,z\) implies
\[
(q-F)-z(q-1)=1-F-(z-1)(q-1)\le0.
\]
Since \(\kappa+w\ge2\), the alternative \(C\le-1\) would make
\(N_t\le-2\). Hence \(C\ge0\), and the unchanged mask argument gives
the bounds (2.9).

Put \(Z=|Y|\). The odd substitution excludes \(s=0\), so
\[
\begin{gathered}
Z\ge q^3,\quad X\ge0,\quad |a|=(X+1)Z,\quad
k(Z-1)<|c|<2kZ,\quad |H|>1,\\
\Delta\ge(q^3-1)(q^3-3)>q^4>R.
\end{gathered}
\tag{6.5}
\]
Since \(i\ne0\) and \(|c|>1\),
\[
D>\Delta+|c|,\qquad
1<b=c+RD<Q.
\tag{6.6}
\]
For the upper bound, explicitly,
\[
Q-b=(\Delta-R)D-\Delta-c
 \ge D-\Delta-|c|>0.
\]
Also \(G,D,Q>1\) and
\(\ell=XkY^2\ge0\), so \(A\ge1\).
In particular \(\mathcal B>0\), even with signed or zero \(T\).

\subsection{Root recovery with possibly negative weights}

The factorization of the norm supplies a coherent triple
\(u^2=G,\ t^2=A,\ f^2=D\) and
\[
y=\frac{W-Lu-HLt}{H},\qquad
dy^2=e-B_af=Q(vf-b)^2-1.
\tag{6.7}
\]
Work in a finite normal splitting field. Every conjugate \(y'\)
satisfies \(|y'|<\mathcal B\).
To see this, put \(M_0=|v|\sqrt D+b>1\).
The real quantity \(Q(vf'-b)^2-1\) lies between \(-1\) and
\(QM_0^2-1\), so its absolute value is at most \(QM_0^2-1\).
Thus
\[
|y'|^2\le\frac{QM_0^2-1}{Q-1}<2M_0^2,\qquad
\mathcal B=v^2+D+2b\ge2M_0.
\tag{6.8}
\]
This bounds imaginary as well as real branches.

Set \(U_0=|H|\mathcal B>2\).
Although \(L\) can be negative, \(|L|\ge U_0-1\).
Since \(\sqrt G>2\), we obtain
\[
|L|\sqrt G>U_0.
\tag{6.9}
\]
If \(X>0\), then \(X\ge q\), and
\[
\sqrt A>\ell=kZ|E|,\qquad
\sqrt G<(|a|+2)|c|<2(|a|+2)kZ.
\]
Consequently
\[
\begin{split}
|H|\sqrt A-\sqrt G
&>kZ\bigl((4|a|-3)|E|-2(|a|+2)\bigr)\\
&\ge kZ(2|a|-7)>2.
\end{split}
\]
Using \(2|L|>U_0\), this implies the second strict separation
\[
|HL|\sqrt A>|L|\sqrt G+|H|\mathcal B.
\tag{6.10}
\]
If \(t\) were irrational, an automorphism of the splitting field
would flip its sign. Subtracting the conjugates of
\(W=Lu+HLt+Hy\) would contradict (6.10), because the other two
differences have absolute values at most \(2|L|\sqrt G\)
and strictly less than \(2|H|\mathcal B\).
Thus \(t\) is rational.
If \(X=0\), then \(A=1\), so \(t\) is rational directly
and (6.10) is unnecessary.

With \(t\) rational, an automorphism flipping \(u\) would contradict
(6.9). Hence \(u\) is rational as well. Both are integers because
their squares are integers; equation (6.7) then makes \(y\) rational.

\subsection{Saturation and the two exceptional boundaries}

If \(T\ne0\), then \(v\ne0\), \(B_a\ne0\), and
\[
f=\frac{e-dy^2}{B_a}\in\mathbb Q.
\]
Since \(f^2=D\), this root is integral. Now \(vf-b\in\Z\), so
Lemma~\ref{lem:saturation} applied with
\(m=|\Delta ic^2|\) makes \(y\) integral.
If \(T=0\), the same lemma applies directly with \(V=b\),
without first recovering \(f\), and again gives \(y\in\Z\).
Because \(L\equiv1\pmod H\), in either case
\[
\gamma=\frac{u-b_0}{H}
=r-\mathcal B u-Lt-y\in\Z.
\tag{6.11}
\]

If \(X=0\), then \(u=ac+H\gamma\).
Reducing \(u^2-\Delta c^2=1\) modulo \(H\), with
\(\Delta=a^2+H\), gives \(H\mid1\), contrary to \(|H|>1\).
Therefore \(X>0\) and \(w>0\).
If \(T=0\), (6.7) gives
\[
y^2=b^2+\frac{b^2-1}{Q-1}.
\]
The integer inequalities \(1<b<Q\) put this number strictly
between \(b^2\) and \((b+1)^2\), a contradiction.
Thus \(T\ne0\), and \(f\) is integral in all possible cases.

We have now reconstructed all four equations (2.8), with signed
\(t,u,f,y,\gamma,h,T\). No division by an unproved nonzero
\(T\) or \(B_a\) was used.

\subsection{Why a negative odd scale is impossible}

The recovered equations meet the hypotheses of the signed-scale
exclusion in package 23 \cite{oddscale}. The normalized multiplier
need only be nonzero: replacing \(i\) by \(|i|\) changes none of
the equations.
For clarity, the mechanism of that exclusion is as follows.
If \(s<0\), put
\[
Z=-Y,\qquad B_0=(X+1)Z-2,\qquad c_0=-c,\qquad
P_0=2XZ^2+1.
\]
The first and main Pell classifications give indices \(n,p\)
with \(k=2\psi_{P_0}(n)\) and \(c_0=\psi_{B_0}(p)\).
Here \(\psi_B\) denotes the ordinate of
\((B+\sqrt{B^2-1})^j\).
The normalized strong equation makes its index divisible by
\(pc_0\). The auxiliary congruences then recover \(p=R\);
the sign of \(c\) disappears in the squared congruence.
The strict index interval forces \(\epsilon=1\) and
\(R=2r_0+1\), and the Pell growth comparison gives
\[
Z>\frac{X^{r_0}}2>2^R.
\]
The two possible signs of the main root give respectively
\[
X+2^R\equiv0\pmod H
\quad\text{or}\quad
X\,2^R\equiv1\pmod H.
\]
Both are excluded by
\[
0<X+2^R<|H|,
\qquad
1<X\,2^R<XZ<|H|.
\]
The full inherited rank argument, including its signed congruences,
is supplied in the reference copy of
\texttt{ODD\_SCALE\_27\_PROOF.md}.
It applies to the reconstructed core, not to a merely rational
norm witness. This order is essential.

Since \(s\) is odd and cannot be negative, it is positive.
The signed normalized compiler now decodes a blank \(\KTM\) halt
and forces \(N_k=N_t=1\).

\subsection{Simultaneous completeness of the domain restrictions}

From an actual blank halt, canonical independent width and height
padding gives \(q=2^\theta\) with \(\theta=d_cN\) odd.
The canonical rank \(R\) is odd, so
\[
w=2^{R-\theta}
\]
is a positive square.
Choosing the normalized strong index \(m=cR\) gives odd positive
\(i\), together with the required auxiliary quotient, as in
the signed-domain proof \cite{signed}.
These restrictions concern the same padded tableau.

To check odd \(s\) without changing that witness, set
\[
S=(C-1)+qF,\qquad
T_{\rm mask}=(M_CJ+1)+q(M_FJ-1),
\tag{6.12}
\]
where \(M_F\) is the \emph{native} high mask, not the shifted
circuit port \(\widehat M_F\).
The native mask constraints give
\(S\mathbin{\&}T_{\rm mask}=0\), both quantities below \(q^2\),
and \(\pc(T_{\rm mask})=\theta+2\).
The exact packing identity is
\[
R=(q^2-1-S)q^2+(S+T_{\rm mask}).
\]
The two blocks are disjoint, so
\[
\pc(R)=2\theta+\pc(T_{\rm mask})=3\theta+2.
\]
The inherited half-binomial valuation gives
\(v_2(Y)=\pc(R)-2=3\theta\). Therefore \(s=Y/q^3\) is positive odd.
Every remaining guarded coordinate is in the image of (6.2).
Finally, with \(t\) again the canonical first Pell root, choose
\[
r=\gamma+\mathcal B u+Lt+y.
\tag{6.13}
\]
This integer makes one norm factor vanish, proving completeness
for (6.4) with all three restrictions satisfied simultaneously.

\subsection{The system and sum-of-squares alternatives}

The unshifted system is \(N_k=1,\ N_t=1,\ \mathcal N=0\).
For the recommended shorter system, use six positive maps including
the new \(z\), put \(z_{\rm old}=z+1\), and write
\[
\boxed{
\begin{aligned}
N_k&=1,\\
(\kappa+w)C&=F+z(q-1),\\
N_0^2&=GN_1^2.
\end{aligned}}
\tag{6.14}
\]
The proof of the shift is the same as for (5.14), after the core
has been decoded. This system has \(24\) integer unknowns,
three equations, \(157\) gates, and residual degrees \(12,4,2416\).
The unshifted version uses \(159\) gates.
For a single polynomial obtained directly from the system, the
archive also includes
\[
(N_k-1)^2+(N_t-1)^2+\mathcal N^2=0,
\tag{6.15}
\]
with degree \(4832\) and \(166\) operations.
Its integer zero set is exactly that of the unshifted system.

\subsection{Operation count and exact degree}

Each positive map costs seven operations, the two odd maps cost two
each, and the square map costs one. The total domain-conversion
cost is therefore \(42+4+1=47\).
The \(116\)-gate package-21 product circuit is retained literally
after that prefix, yielding \(163\) gates.
This is a new theorem about the domain of an old polynomial
circuit, as well as a smaller explicit integer presentation.

Give \(J,F,\alpha,z,w,\eta,\zeta\) weight two and the other
core coordinates weight one. Key weighted degrees are
\[
\begin{array}{c|rrrrrrrr}
 &c&H&\Delta&A,G&D&Q&b,\mathcal B&L\\ \hline
\deg&9&11&22&40&60&82&68&79
\end{array}
\]
In the evaluator (3.4), \(n\) has weight \(220\), and the unique
highest part of \(C_0\) is \(dn\), of weight \(302\).
The highest part of \(N_0\) has weight \(1208\), while
\(\deg N_1\le1087\).
Thus \(N_0^2\) has weight \(2416\) and
\(GN_1^2\) has weight at most \(2214\); there is no leading
cancellation. With \(q_0=\beta J,\ k_0=\eta+\zeta\), the
highest weighted part of the norm is
\[
2^{80}i^{48}k_0^{112}w^{112}s^{224}q_0^{832}(q_0-F)^{16}.
\tag{6.16}
\]
Multiplying by the leading outer factors gives
\[
\begin{split}
(P_{24})_{\rm top}={}&-2^{80}h i^{48}k_0^{112}w^{113}s^{225}
 q_0^{836}(q_0-F)^{16}\\
&\hspace{9mm}\cdot[w(q_0-F-\alpha)-zq_0],
\end{split}
\tag{6.17}
\]
before the actual degree-two substitutions.
For a guarded coordinate \(p\), denote the first variable in its
positive-map triple by \(p_x\).
After substituting, the coefficient of
\[
J_x^{1706}h z_x^2 i_0^{48}\eta_x^{224}w_0^{226}s_0^{225}
\]
is \(2^{353}\beta^{853}\ne0\).
Its degree is \(2432\).
The analogous norm coefficient is
\(2^{352}\beta^{848}\ne0\), at degree \(2416\).
Consequently the product and system degrees hold for the actual
fixed coefficients, uniformly in every authentic \(\beta>0\).
